%%%%%%%%%%%%%%%%%%%%%%%%%%%%%%%%%%%%%%%%%%%%%%%%%%%%%%%%%%%%%%%%%%%%%%%%%%%%%%%%%
%	ABSTRACT PAGE
%   TICKET-23 --- Abstract (~250 words)
%%%%%%%%%%%%%%%%%%%%%%%%%%%%%%%%%%%%%%%%%%%%%%%%%%%%%%%%%%%%%%%%%%%%%%%%%%%%%%%%%
\addtotoc{Abstract}
\abstract{\addtocontents{toc}{\vspace{1em}}

The Traveling Salesman Problem (TSP) is a classical combinatorial
optimisation problem that seeks the shortest closed tour visiting each city
exactly once.  Because the search space grows factorially with the number of
cities, exact methods become impractical for non-trivial instances, making
metaheuristic approaches such as Genetic Algorithms (GAs) attractive
alternatives.  A persistent challenge when applying GAs to the TSP is that
standard crossover and mutation operators can produce infeasible solutions
--- chromosomes that visit some cities more than once while omitting others
--- leading to meaningless fitness evaluations and distorted selection
pressure.

This report investigates the design, implementation, and critical analysis
of a GA for the static single-objective TSP with a focus on feasibility
repair mechanisms.  We implement three permutation-preserving crossover
operators (PMX, OX, CX), two mutation operators (swap, inversion), and a
repair mechanism that diagnoses constraint violations and restores
feasibility through duplicate removal and missing-city insertion.  Using the
kroA100 benchmark instance (100 cities, known optimal tour length 21\,282),
we conduct a comparative experiment of 30 independent runs per
constraint-handling strategy, evaluating solution quality, convergence
behaviour, and population diversity.  A Wilcoxon signed-rank test confirms a
statistically significant difference between repair-enabled and unrepaired
runs ($p < 0.001$), demonstrating that without repair, the GA converges to
artificially low but invalid fitness values.  Parameter sensitivity analysis
further explores the interaction between GA configuration and the repair
mechanism's effectiveness.  Our findings highlight the critical role of
constraint handling in permutation-based evolutionary optimisation and
provide empirical evidence for the trade-offs between feasibility
enforcement and search-space exploration.
}
\clearpage
